\section{绪论}

\subsection{研究背景}

视觉语言模型在近几年发展迅速。CLIP 等模型通过海量图文数据的对比训练，实现了图像和文本在统一语义空间中的对齐；BLIP-2、LLaVA 和 Qwen-VL 等生成式模型进一步将视觉编码器与大语言模型连接，具备了图像描述、视觉问答等能力。然而，这些模型对输入图像的质量高度敏感——在低分辨率、模糊、噪声或遮挡条件下，语义识别和描述生成的性能会明显下降。

与此同时，脑电图作为一种成熟的脑机接口技术，能以毫秒级时间分辨率记录视觉刺激引发的神经电活动。我在查阅相关文献时了解到，已有研究者尝试从 EEG 中解码视觉类别、检索语义相关图像，甚至将 EEG embedding 映射到 CLIP 空间。这让我产生了一个想法：当视觉输入质量受损时，人脑的神经响应是否能作为额外的信息来源，帮助机器视觉系统弥补退化带来的信息损失？带着这个问题，我开始了本次课程设计的探索。

\subsection{课程设计目标与意义}

本次课程设计的目标是在 Thought2Text 小样本 EEG-image 数据集上，构建并评估一个视觉退化条件下的 EEG 辅助视觉语言理解系统。具体来说，我希望通过这个项目完成以下事情：设计 EEG 编码器从脑电信号中提取有用的语义特征；将 EEG 特征与退化图像的视觉特征进行融合；对比真实配对 EEG、打乱 EEG 和随机噪声 EEG，客观判断 EEG 是否提供了有效的语义辅助；最终尝试将融合后的特征输入生成式 VLM，产出自然语言图像描述。

从个人学习和实践的角度来看，这个课题让我有机会将课堂上学的深度学习、Transformer、多模态融合等知识点串起来，在一个真实的数据集上完整体验从问题定义、数据处理、模型设计、训练调参到结果评估的全流程。特别是在单 GPU 条件下完成大模型微调、处理 EEG 这种非标准模态数据，都是很有价值的工程实践经历。

\subsection{研究内容与技术路线}

本次课程设计主要完成了四项工作：

\begin{enumerate}
  \item 搭建视觉退化条件下 EEG 辅助语义识别和生成式描述的完整实验框架，定义了六种退化条件和四种 EEG 对照模式；
  \item 设计并实现 A2 temporal-spectral-spatial EEG 编码器，从 EEG 信号的时间、频谱和空间通道三个维度并行提取特征，通过门控残差融合对退化图像进行语义补充；
  \item 实现 VTF token-level 融合方法，将 EEG 信息注入到 CLIP ViT 的 visual patch token 中，为后续接入生成式模型做准备；
  \item 构建基于 Qwen2-VL 的生成式图像描述模块，通过 LoRA 在单 GPU 上高效微调 2B 模型，并对比了五种生成技术方案的优劣。
\end{enumerate}

在技术路线上，我采取了由简到繁的分步推进策略：先用 CLIP 的预训练语义空间作为跨模态桥梁（降低对 EEG 数据量的依赖），再在分类任务上验证 EEG 的效果（可量化、可对照），然后逐步推进到 token-level fusion 和生成式输出。每一步都有明确的中间检查点，出错时方便定位问题。

\subsection{报告结构安排}

本报告分为十章。第 1 章绪论；第 2 章相关技术基础；第 3 章数据集与任务；第 4 章系统总体设计；第 5 章 A2 EEG 语义融合模型；第 6 章 token-level EEG 视觉增强；第 7 章生成式 VLM 设计；第 8 章实验结果与分析；第 9 章工程实现与代码；第 10 章总结与反思。
